\onecolumn
\section{Evidence and diagnostics annex}
\small

\subsection{Evidence tiers and provenance}

Tier A evidence uses standardized five-seed rollout tables on the registered 61 by 41 grid. Tier B evidence uses frozen checkpoints and registered diagnostic state sets after method selection. Tier C evidence contains short screens, local parameter perturbations, and one-seed probes. Only Tier A results determine the headline methods. Tier B evidence supports mechanism interpretation. Tier C evidence is exploratory.

The selected mixed actor family is the registered no-shift automatic-priority DAgger follow-up. The selected pure deployment is the registered SimbaV2 symmetric-actor Q41 unanimous-search evaluation. Both contain five actor indices and 12,505 fixed-grid rollouts.

The mixed initializer command used 400,000 static labels, three learner-controlled DAgger rounds, and a 20\% near-upright sampler within 35 degrees. The follow-up used 240,000 reset-support labels plus 20,000 automatically prioritized learner states. This provenance is reported as a sampler limitation. The deployment actor and local Q-search do not read the reference or use an angle rule.

Demonstration-augmented policy gradients provide another route for combining imitation and reward optimization \cite{rajeswaran2018dapg}. The selected mixed route instead keeps the supervised actor target separate from the reward-trained deployment critic.

\subsection{Systematic two-repository census}

\begin{center}
\includegraphics[width=.96\textwidth]{27_forensic_audit_summary.png}
\captionof{figure}{Forensic inventory of standardized evaluations in Project 15 and Project 15-sac-n-experiments. Counts describe the crawl before the final project decision retained the accepted mixed result with a sampler footnote. The result tables in the main report follow the final decision.}
\end{center}

The inventory is useful for identifying missing seeds, aliases, deployment variants, and unreported experiment families. Its earlier binary training-sampler classification is not used to replace the accepted mainline result.

\clearpage
\subsection{FastSACN8 target weight and learning curves}

\begin{center}
\includegraphics[width=.96\textwidth]{36_fastsacn_objective_share.png}
\captionof{figure}{Normalized loss coefficients for FastSACN target variants. The $\lambda=0.5$ FastSACN8 recipe assigns 0.775\% of the normalized coefficient to the eight-step categorical cross-entropy term. This is not a measured share of the gradient or cumulative learning effect.}
\end{center}

\begin{center}
\includegraphics[width=.96\textwidth]{14_fastsacn8_loss_curves.png}
\captionof{figure}{Stored loss and evaluation traces for the earlier FastSACN8 family. Loss scales differ across objectives, so the curves describe optimization history rather than provide a direct cross-objective ranking.}
\end{center}

The historical FastSACN8 result changes training length, critic update dose, and deployment relative to the selected one-step pure method. The completed five-seed combination keeps the 100k budget, SimbaV2 architecture, and deployment fixed while changing both the target family and critic update dose. On 27,765 locked off-grid trials per arm, its selected deployment improves near-reference reliability from 94.489\% to 95.909\% and strict wins from 17.792\% to 21.930\%, while task success decreases from 93.038\% to 92.948\%. The authority grid records 12,008 of 12,505 near-reference successes, 11,614 task successes, and 2,854 strict wins. This comparison tests the combined target-and-dose recipe rather than the target family alone.

\begin{center}
\includegraphics[width=.96\textwidth]{47_actor_critic_transfer.png}
\captionof{figure}{Five-seed target ablations and the selected pure-RL deployment intervention. FastSACN8 changes both the target and critic update dose. Action search then uses critic preferences to replace many actions chosen by the actor.}
\end{center}

The seed-zero actor-UTD2 screen improves strict wins and mean return relative to the selected seed-zero configuration, but lowers near-reference reliability from 96.668\% to 96.128\% and task success from 92.527\% to 92.112\%. It was not extended to four more seeds and is not a mainline result.

\clearpage
\subsection{Actor geometry and saturation}

\begin{center}
\includegraphics[width=.96\textwidth]{16_completed_target_actor_geometry.png}
\captionof{figure}{Checkpoint-resolved actor saturation, action sensitivity, and reflection error for completed temporal-target arms. These actor-geometry measures are descriptive and do not identify the effect of critic update dose.}
\end{center}

Relative to the five one-step actors, the P7 combination raises the median saturation fraction from 59.1\% to 70.3\% and lowers the median tanh action-sensitivity statistic from 0.120 to 0.089. Both paired changes have the same direction in all five indexed seeds. Its ordinary actor is less reliable, while its critic-guided deployment is more reliable. This pattern associates the gain with deployment repair rather than a better raw actor, but the target and critic dose change together.

\begin{center}
\includegraphics[width=.94\textwidth]{13_fastsacn8_seed0_geometry.png}
\captionof{figure}{Seed-0 FastSACN8 actor geometry probe. This one-seed analysis is exploratory and cannot become a headline result.}
\end{center}

Pure-actor failures are strongly associated with saturated actions, but the reference action is also saturated on nearly all of those states. Saturation therefore identifies a difficult swing-up regime rather than proving that the actor chose the wrong direction. The first-action intervention separates a smaller opposite-sign error group from a larger same-sign timing group.

\clearpage
\subsection{Dormancy, effective rank, and telemetry}

\begin{center}
\includegraphics[width=.97\textwidth]{24_all_arm_training_telemetry.png}
\captionof{figure}{Seed ranges for online evaluation, actor loss, critic loss, maximum measured dormancy, and minimum effective rank across completed five-seed pure-RL arms.}
\end{center}

The separation between generated experience and the policy extracted from it follows exploration-versus-optimization diagnostics \cite{berseth2025exploration}. The registered threshold follows dormant-neuron diagnostics \cite{sokar2023redo} and finds zero dormant units in the completed SimbaV2 arms. This is a negative result. It rules out widespread measured dormancy as the explanation for the selected reliability gap. It does not establish healthy plasticity, well-conditioned gradients, or good critic calibration. Effective rank, primacy bias, plasticity, and policy churn are narrower complementary checks \cite{gulcehre2022rank,nikishin2022primacy,lyle2023plasticity,nauman2024plasticity,tang2024churn}. Loss magnitudes also differ across target objectives and should not be ranked without normalization.

\subsection{Local parameter landscapes}

\begin{center}
\includegraphics[width=.94\textwidth]{12_parameter_landscape_summary.png}
\captionof{figure}{Local parameter perturbation planes around selected frozen checkpoints. The geometry depends on perturbation directions, parameter scaling, and the fixed diagnostic batch.}
\end{center}

The planes show local sharpness and coupling for the chosen directions. They do not provide a global loss landscape or a causal explanation. Their role is to check whether a method ends in an obviously unstable local neighborhood. The current result is exploratory because the selected categories use different objectives and datasets.

\clearpage
\subsection{First disagreement and critic ranking}

\begin{center}
\includegraphics[width=.99\textwidth]{22_first_disagreement_survival.png}
\captionof{figure}{Step of first large actor-reference action disagreement on a locked off-grid set. Thin curves show indexed actor seeds and thick curves pool failure rows. Different failure sets make the category comparison descriptive.}
\end{center}

\begin{center}
\includegraphics[width=.96\textwidth]{28_critic_ranking_calibration.png}
\captionof{figure}{Critic proposal calibration at first-disagreement states. Precision, recall, balanced accuracy, AUC, and Q-margin bins separate useful local ranking from harmful proposals.}
\end{center}

For pure RL, the pooled first-disagreement diagnostic contains 2,259 failure rows. Helpful-proposal precision is 0.782, recall is 0.703, balanced accuracy is 0.623, and ROC AUC is 0.655. The critic contains useful information but does not reliably order every hard-state proposal. This result is consistent with the positive net gain and nonzero break count of global Q-search.

\clearpage
\subsection{Seed-resolved timing}

\begin{center}
\includegraphics[width=.99\textwidth]{23_actor_timing_baselines.png}
\captionof{figure}{Seed-resolved trajectory timing for pure RL, the selected mixed pipeline, and DAgger timing controls. Counts are displayed because the failure populations differ substantially.}
\end{center}

The mixed failures diverge later than pure failures on the selected frozen sets. This observation is consistent with learner-state supervision reducing compounding drift. It is not a causal estimate because the pipelines differ in labels, initialization, interaction counts, critics, and selection history.

\subsection{Closed-loop sensitivity}

\begin{center}
\includegraphics[width=.94\textwidth]{15_closed_loop_sensitivity.png}
\captionof{figure}{Closed-loop response to small state perturbations on selected frozen policies. The diagnostic measures trajectory amplification after selection.}
\end{center}

Closed-loop amplification tests whether small early deviations grow into a reliability failure under the same frozen actor. It complements the action-ranking diagnostic: a good local action is insufficient when later policy actions amplify residual state error.

\clearpage
\subsection{Reference-prefix specificity}

\begin{center}
\includegraphics[width=.95\textwidth]{40_p7_reference_prefix_intervention.png}
\captionof{figure}{Exact-lineage repair rate as a state-matched reference policy controls an initial prefix and then returns control to each selected P7 raw actor.}
\end{center}

\begin{center}
\includegraphics[width=.92\textwidth]{41_p7_reference_prefix_specificity.png}
\captionof{figure}{Exact-lineage specificity controls. Shuffled reference actions and zero torque do not reproduce state-matched repair at longer prefixes.}
\end{center}

The intervention demonstrates sufficiency on the tested frozen failures. A sustained state-matched prefix repairs nearly every trial, while generic interruptions do not. The reference is diagnostic only and is unavailable at deployment. This result does not show which training update would teach the same corrective sequence.

\subsection{Interpretation boundary}

The combined diagnostics reject several simple explanations at their measured resolution:

\begin{itemize}
\item Failed seeds do not systematically have farther nearest training resets.
\item Reference-like actions are usually present in nearby replay.
\item Complete stored trajectories rarely outperform the terminal actor by more than five return units.
\item Widespread dormant units are absent at the registered threshold.
\end{itemize}

They support a narrower account:

\begin{itemize}
\item Pure failures often begin with early actor-reference disagreement.
\item One local action rarely removes the complete return deficit.
\item State-matched corrective sequences are sufficient on almost all tested failures.
\item The learned critic ranks helpful proposals better than chance but makes consequential tail errors.
\end{itemize}

These conclusions apply to the selected Pendulum checkpoints and diagnostic state sets. Broader claims require other environments, matched starts and budgets, independent critics, and a uniform initializer.
